\documentclass[12pt]{article}

% ============================================================
% Packages
% ============================================================

\usepackage{geometry}
\usepackage{titling}
\usepackage{amsmath,amssymb}
\usepackage{fontspec}
\usepackage{xcolor}
\usepackage{listings}
\usepackage{minted}
\usepackage{tikz}
\usepackage{nicematrix}
\usepackage{unicode-math}

% ============================================================
% Page layout and fonts
% ============================================================

\geometry{
    a4paper,
    top=2.5cm,
    bottom=2.5cm,
    left=2.5cm,
    right=2.5cm
}

\setlength{\droptitle}{-3cm}

\setmainfont{Libertinus Serif}
\setsansfont{Libertinus Sans}
\setmonofont{DejaVu Sans Mono}
\setmathfont{STIXTwoMath-Regular.otf}

\usetikzlibrary{fit}

% ============================================================
% Custom commands
% ============================================================

\newcommand{\Matrix}[1]{\mathbf{#1}}
\newcommand{\Vector}[1]{\mathbf{#1}}
\newcommand{\Scalar}[1]{#1}

\newcommand{\vsp}{\vspace*{1em}}
\newcommand{\snl}{\vspace*{2pt}}

% ============================================================
% Title
% ============================================================

\title{\textbf{Affine Functions}}
\author{Devi Prasad M}
\date{August 2026}

% ============================================================
% Document
% ============================================================

\begin{document}

\maketitle

% ============================================================
\section{Affine Functions and Affine Combinations}
% ============================================================

An \emph{affine function} is a linear function followed by a
translation. It can be written as

\[
    f(\Vector{v}) = L(\Vector{v}) + \Vector{b},
\]

where

\begin{itemize}
    \item $L$ is a linear function, and
    \item $\Vector{b}$ is a fixed translation vector.
\end{itemize}

The linearity of $L$ means that

\[
L\left(
    \sum_{i=1}^{k} c_i\Vector{v}_i
\right)
=
\sum_{i=1}^{k} c_iL(\Vector{v}_i).
\]

A natural question is:

\begin{quote}
\emph{Does an affine function preserve linear combinations in the
same way that a linear function does?}
\end{quote}

The answer is almost the same, but with one important restriction:

\[
\boxed{
\text{An affine function preserves affine combinations.}
}
\]

An \emph{affine combination} is a weighted combination whose
coefficients sum to $1$.

% ============================================================
\subsection{Why Must the Coefficients Sum to 1?}
% ============================================================

Consider a weighted combination of vectors

\[
    \Vector{v}
    =
    \sum_{i=1}^{k} c_i\Vector{v}_i.
\]

Applying the affine function gives

\begin{align*}
f\left(
    \sum_{i=1}^{k} c_i\Vector{v}_i
\right)
&=
L\left(
    \sum_{i=1}^{k} c_i\Vector{v}_i
\right)
+
\Vector{b}
\\[4pt]
&=
\sum_{i=1}^{k} c_iL(\Vector{v}_i)
+
\Vector{b}.
\end{align*}

Now consider the same weighted combination of the outputs:

\begin{align*}
\sum_{i=1}^{k} c_i f(\Vector{v}_i)
&=
\sum_{i=1}^{k}
c_i\left(
    L(\Vector{v}_i)+\Vector{b}
\right)
\\[4pt]
&=
\sum_{i=1}^{k}c_iL(\Vector{v}_i)
+
\left(
    \sum_{i=1}^{k}c_i
\right)\Vector{b}.
\end{align*}

For the two expressions to be equal, we require

\[
\Vector{b}
=
\left(
    \sum_{i=1}^{k}c_i
\right)\Vector{b}.
\]

For an arbitrary translation vector $\Vector{b}$, this requires

\[
\boxed{
\sum_{i=1}^{k}c_i=1.
}
\]

Therefore,

\[
\boxed{
f\left(
    \sum_{i=1}^{k}c_i\Vector{v}_i
\right)
=
\sum_{i=1}^{k}c_i f(\Vector{v}_i)
\qquad
\text{whenever}
\qquad
\sum_{i=1}^{k}c_i=1.
}
\]

This is the fundamental preservation property of an affine function.

% ============================================================
\subsection{Linear Combination vs.\ Affine Combination}
% ============================================================

It is important to distinguish between a linear combination and an
affine combination.

A \emph{linear combination} has the form

\[
    \sum_{i=1}^{k}c_i\Vector{v}_i,
\]

with no restriction on the coefficients.

An \emph{affine combination} has the additional constraint

\[
    \boxed{
    \sum_{i=1}^{k}c_i=1.
    }
\]

The coefficients may be positive, zero, or negative.

For example,

\[
    0.4\Vector{v}_1+0.6\Vector{v}_2
\]

is an affine combination because

\[
    0.4+0.6=1.
\]

On the other hand,

\[
    2\Vector{v}_1+3\Vector{v}_2
\]

is a linear combination, but not an affine combination, because

\[
    2+3=5\neq1.
\]

If, in addition,

\[
    c_i\geq0
    \qquad\text{for all }i,
\]

then the affine combination is called a
\emph{convex combination}.

For example,

\[
    0.4\Vector{v}_1+0.6\Vector{v}_2
\]

is both an affine combination and a convex combination.

% ============================================================
\section{A Concrete 2D Example}
% ============================================================

Consider the affine function

\[
    f(\Vector{v})
    =
    \Matrix{L}\Vector{v}+\Vector{b},
\]

where

\[
\Matrix{L}
=
\begin{bmatrix}
2 & 0\\
0 & 3
\end{bmatrix},
\qquad
\Vector{b}
=
\begin{bmatrix}
4\\
-1
\end{bmatrix}.
\]

The matrix $\Matrix{L}$ doubles the first coordinate and triples
the second coordinate.

The vector $\Vector{b}$ then translates the result by

\[
4 \text{ units in the }x\text{-direction}
\]

and

\[
-1 \text{ unit in the }y\text{-direction}.
\]

Let

\[
\Vector{v}_1
=
\begin{bmatrix}
1\\
2
\end{bmatrix},
\qquad
\Vector{v}_2
=
\begin{bmatrix}
3\\
6
\end{bmatrix}.
\]

We choose the coefficients

\[
    c_1=0.4,
    \qquad
    c_2=0.6.
\]

Notice that

\[
    c_1+c_2=0.4+0.6=1.
\]

Therefore,

\[
    0.4\Vector{v}_1+0.6\Vector{v}_2
\]

is an affine combination.

Geometrically, it lies on the line segment joining
$\Vector{v}_1$ and $\Vector{v}_2$.

% ============================================================
\subsection{Method 1: Combine the Inputs First}
% ============================================================

First construct the affine combination:

\begin{align*}
\Vector{v}_3
&=
0.4\Vector{v}_1+0.6\Vector{v}_2
\\[4pt]
&=
0.4
\begin{bmatrix}
1\\
2
\end{bmatrix}
+
0.6
\begin{bmatrix}
3\\
6
\end{bmatrix}
\\[6pt]
&=
\begin{bmatrix}
0.4\\
0.8
\end{bmatrix}
+
\begin{bmatrix}
1.8\\
3.6
\end{bmatrix}
\\[6pt]
&=
\begin{bmatrix}
2.2\\
4.4
\end{bmatrix}.
\end{align*}

Now apply the affine function:

\begin{align*}
f(\Vector{v}_3)
&=
\begin{bmatrix}
2 & 0\\
0 & 3
\end{bmatrix}
\begin{bmatrix}
2.2\\
4.4
\end{bmatrix}
+
\begin{bmatrix}
4\\
-1
\end{bmatrix}
\\[6pt]
&=
\begin{bmatrix}
4.4\\
13.2
\end{bmatrix}
+
\begin{bmatrix}
4\\
-1
\end{bmatrix}
\\[6pt]
&=
\boxed{
\begin{bmatrix}
8.4\\
12.2
\end{bmatrix}
}.
\end{align*}

% ============================================================
\subsection{Method 2: Transform the Inputs First}
% ============================================================

Now let us reverse the order.

First transform $\Vector{v}_1$:

\begin{align*}
f(\Vector{v}_1)
&=
\begin{bmatrix}
2 & 0\\
0 & 3
\end{bmatrix}
\begin{bmatrix}
1\\
2
\end{bmatrix}
+
\begin{bmatrix}
4\\
-1
\end{bmatrix}
\\[6pt]
&=
\begin{bmatrix}
2\\
6
\end{bmatrix}
+
\begin{bmatrix}
4\\
-1
\end{bmatrix}
\\[6pt]
&=
\begin{bmatrix}
6\\
5
\end{bmatrix}.
\end{align*}

Similarly,

\begin{align*}
f(\Vector{v}_2)
&=
\begin{bmatrix}
2 & 0\\
0 & 3
\end{bmatrix}
\begin{bmatrix}
3\\
6
\end{bmatrix}
+
\begin{bmatrix}
4\\
-1
\end{bmatrix}
\\[6pt]
&=
\begin{bmatrix}
6\\
18
\end{bmatrix}
+
\begin{bmatrix}
4\\
-1
\end{bmatrix}
\\[6pt]
&=
\begin{bmatrix}
10\\
17
\end{bmatrix}.
\end{align*}

Now form the same affine combination of the outputs:

\begin{align*}
0.4f(\Vector{v}_1)+0.6f(\Vector{v}_2)
&=
0.4
\begin{bmatrix}
6\\
5
\end{bmatrix}
+
0.6
\begin{bmatrix}
10\\
17
\end{bmatrix}
\\[6pt]
&=
\begin{bmatrix}
2.4\\
2.0
\end{bmatrix}
+
\begin{bmatrix}
6.0\\
10.2
\end{bmatrix}
\\[6pt]
&=
\boxed{
\begin{bmatrix}
8.4\\
12.2
\end{bmatrix}
}.
\end{align*}

Thus,

\[
\boxed{
f(0.4\Vector{v}_1+0.6\Vector{v}_2)
=
0.4f(\Vector{v}_1)+0.6f(\Vector{v}_2)
}.
\]

The two paths produce exactly the same result.

% ============================================================
\section{Why Does This Work?}
% ============================================================

The equality is not a coincidence arising from our particular
numerical example.

Consider the general case

\[
    f(\Vector{v})=L(\Vector{v})+\Vector{b}.
\]

Then

\begin{align*}
f(c_1\Vector{v}_1+c_2\Vector{v}_2)
&=
L(c_1\Vector{v}_1+c_2\Vector{v}_2)+\Vector{b}
\\
&=
c_1L(\Vector{v}_1)
+
c_2L(\Vector{v}_2)
+
\Vector{b}.
\end{align*}

On the other hand,

\begin{align*}
c_1f(\Vector{v}_1)+c_2f(\Vector{v}_2)
&=
c_1
\left(
L(\Vector{v}_1)+\Vector{b}
\right)
+
c_2
\left(
L(\Vector{v}_2)+\Vector{b}
\right)
\\
&=
c_1L(\Vector{v}_1)
+
c_2L(\Vector{v}_2)
+
(c_1+c_2)\Vector{b}.
\end{align*}

Therefore, the two expressions are equal exactly when

\[
    c_1+c_2=1.
\]

The crucial point is that the translation vector $\Vector{b}$ appears

\[
\textbf{once}
\]

when we first combine the inputs, but it appears

\[
(c_1+c_2)\textbf{ times}
\]

when we first transform the inputs and then combine the outputs.

The condition

\[
    c_1+c_2=1
\]

makes these two quantities identical.


% ============================================================
\section{The Key Idea}
% ============================================================

The difference between linear and affine functions can be summarized
in two equations.

For a linear function,

\[
\boxed{
L\left(
    \sum_i c_i\Vector{v}_i
\right)
=
\sum_i c_iL(\Vector{v}_i)
}
\]

with no restriction on the coefficients.

For an affine function,

\[
\boxed{
f\left(
    \sum_i c_i\Vector{v}_i
\right)
=
\sum_i c_if(\Vector{v}_i)
\qquad
\text{provided}
\qquad
\sum_i c_i=1.
}
\]

The underlying idea is:

\[
\boxed{
\text{Linear maps preserve the origin;
affine maps preserve affine relationships.}
}
\]

The condition

\[
    \boxed{\sum_i c_i=1}
\]

ensures that the translation component is counted exactly once.

\end{document}
